%========================================================
% SECTION 6.2 — Random Phase Ensemble
%========================================================

\newpage

\section{Random Phase Ensemble}

%--------------------------------------------------------
% Subsection 6.2.1 — Objective and Physical Motivation
%--------------------------------------------------------

\subsection*{Objective and Physical Motivation}

The second experiment extends the fixed-phase model by introducing random phase realizations \( \theta_k \), modeling stochastic birefringence effects in realistic optical fibers.

Unlike the coherent regime studied in Section~6.1, the system is no longer described by a single pure state, but by a statistical ensemble of phase-evolved realizations

\[
\left\{
|\psi(\theta_k)\rangle
\right\}_{k=1}^{N},
\]

together with the corresponding ensemble-averaged density operator.

The objective is to characterize how unresolved phase fluctuations modify observable correlations, global coherence, purity, fidelity, and entanglement.

This experiment therefore represents the transition from purely deterministic unitary evolution toward an effective mixed-state description arising from statistical averaging.

%--------------------------------------------------------
% Subsection 6.2.2 — Model and Analytical Predictions
%--------------------------------------------------------

\subsection*{Model and Analytical Predictions}

Each phase realization is described by the pure state

\begin{equation}
|\psi(\theta_k)\rangle
=
\frac{1}{\sqrt{2}}
\left(
|01\rangle + e^{i\theta_k}|10\rangle
\right).
\label{eq:random_phase_realization}
\end{equation}

The ensemble-averaged state is then constructed as

\begin{equation}
\rho^{\mathrm{ens}}
=
\frac{1}{N}
\sum_{k=1}^{N}
|\psi(\theta_k)\rangle
\langle\psi(\theta_k)|.
\label{eq:random_phase_ensemble_density_operator}
\end{equation}

Introducing the ensemble coherence parameter

\begin{equation}
\mu
=
\langle e^{i\theta} \rangle
=
\langle \cos\theta \rangle
+
i
\langle \sin\theta \rangle,
\label{eq:ensemble_coherence_parameter}
\end{equation}

the ensemble-averaged diagonal correlations become

\[
c_{xx}
=
\mathrm{Re}(\mu),
\qquad
c_{yy}
=
\mathrm{Re}(\mu),
\qquad
c_{zz}
=
-1.
\]

The global purity is

\[
\gamma(\rho^{\mathrm{ens}})
=
\frac{1 + |\mu|^2}{2}.
\]

The reduced density operators remain maximally mixed:

\[
\rho_A
=
\rho_B
=
\frac{I}{2},
\]

which implies

\[
\gamma(\rho_A)
=
\gamma(\rho_B)
=
\frac{1}{2}.
\]

The fidelity with respect to the reference Bell state and the concurrence are given by

\[
F_{\Psi^+}
=
\frac{1 + \mathrm{Re}(\mu)}{2},
\qquad
C(\rho^{\mathrm{ens}})
=
|\mu|.
\]

\medskip

The ensemble coherence parameter \( \mu \) therefore provides a compact description of the entire propagation regime. Its real part determines observable transverse correlations and fidelity, while its magnitude controls the residual entanglement of the averaged state.

%--------------------------------------------------------
% Subsection 6.2.3 — Numerical Results
%--------------------------------------------------------

\subsection*{Numerical Results}

Three representative phase distributions are considered:

\begin{itemize}
    \item constant phase distribution;
    \item uniform phase distribution;
    \item Gaussian phase distribution.
\end{itemize}

For each case, the simulator evaluates:

\begin{itemize}
    \item sampled phase statistics;
    \item pure-state correlation trends;
    \item ensemble-averaged correlations;
    \item purity, fidelity, and concurrence.
\end{itemize}

Figure~\ref{fig:experiment2_constant} shows the constant-phase regime.

\begin{figure}[htbp]
\centering
\includegraphics[width=0.9\linewidth]{figures/experiment_2/experiment_2_summary_constant_distribution.png}
\caption{
Constant phase distribution: the ensemble reduces to a coherent realization and concurrence remains maximal.
}
\label{fig:experiment2_constant}
\end{figure}

In this limit, all realizations coincide and the ensemble reproduces the coherent behavior observed in Experiment~1.

\medskip

Figure~\ref{fig:experiment2_uniform} shows the behavior corresponding to a uniform phase distribution.

\begin{figure}[htbp]
\centering
\includegraphics[width=\linewidth]{figures/experiment_2/experiment_2_summary_uniform_distribution.png}
\caption{
Uniform phase distribution: transverse correlations vanish, fidelity decreases, and concurrence tends to zero.
}
\label{fig:experiment2_uniform}
\end{figure}

The averaging process suppresses the transverse coherence terms, leading to vanishing ensemble correlations in the \(x\)--\(y\) plane and to a strong reduction of both fidelity and concurrence.

\medskip

Figure~\ref{fig:experiment2_gaussian} presents the intermediate regime generated by a Gaussian phase distribution.

\begin{figure}[htbp]
\centering
\includegraphics[width=\linewidth]{figures/experiment_2/experiment_2_summary_gaussian_distribution.png}
\caption{
Gaussian phase distribution: partial coherence leads to intermediate values of purity and concurrence.
}
\label{fig:experiment2_gaussian}
\end{figure}

In this case, residual coherence survives ensemble averaging, producing intermediate values of purity, fidelity, and concurrence.

\medskip

The corresponding correlation tensor structures are shown in Figure~\ref{fig:experiment2_tensor_snapshots}.

\begin{figure}[htbp]
\centering

\begin{subfigure}{0.32\linewidth}
    \centering
    \includegraphics[width=\linewidth]{figures/experiment_2/experiment_2_correlation_tensor_constant.png}
    \caption{Constant phase}
\end{subfigure}
\hfill
\begin{subfigure}{0.32\linewidth}
    \centering
    \includegraphics[width=\linewidth]{figures/experiment_2/experiment_2_correlation_tensor_uniform.png}
    \caption{Uniform distribution}
\end{subfigure}
\hfill
\begin{subfigure}{0.32\linewidth}
    \centering
    \includegraphics[width=\linewidth]{figures/experiment_2/experiment_2_correlation_tensor_gaussian.png}
    \caption{Gaussian distribution}
\end{subfigure}

\caption{
Correlation tensor snapshots for different phase distributions.
Random phase averaging suppresses transverse coherence, progressively reducing the off-diagonal tensor components, while the longitudinal correlation \(T_{zz}=-1\) remains invariant.
}
\label{fig:experiment2_tensor_snapshots}
\end{figure}

As the number of realizations increases, the numerically reconstructed tensor components converge toward their analytical expectation values.

At the same time, tensor entries that are analytically zero remain numerically negligible throughout the simulation, confirming the internal consistency of the ensemble implementation.

%--------------------------------------------------------
% Subsection 6.2.4 — Discussion
%--------------------------------------------------------

\subsection*{Discussion}

The numerical results reveal a progressive transition from coherent evolution to ensemble-induced decoherence.

The constant-phase regime reproduces the behavior observed in Experiment~1, where the ensemble effectively reduces to a single pure realization.

As the phase distribution broadens, the transverse coherence terms are progressively suppressed. Consequently, the ensemble coherence parameter satisfies

\[
|\mu|
\rightarrow
0,
\]

leading to the disappearance of coherent transverse correlations.

This transition is reflected directly in the state-characterization quantities:

\begin{itemize}
    \item the global purity decreases from unity;
    \item the reduced purities remain fixed at \(1/2\);
    \item the fidelity decreases;
    \item the concurrence decreases continuously from \(1\) toward \(0\).
\end{itemize}

Unlike the fixed-phase evolution studied previously, the reduction of fidelity is now associated with a genuine degradation of entanglement.

The experiment therefore establishes an important conceptual distinction:

\begin{itemize}
    \item fixed-phase evolution modifies observable correlations without destroying entanglement;
    \item unresolved statistical phase fluctuations generate effective mixedness and progressively suppress entanglement.
\end{itemize}

This mechanism constitutes the first non-unitary effect emerging naturally from the propagation model.

%--------------------------------------------------------
% Subsection 6.2.5 — Physical Interpretation
%--------------------------------------------------------

\subsection*{Physical Interpretation}

Each individual realization of the ensemble remains a pure maximally entangled state. The loss of coherence arises only after averaging over unresolved phase fluctuations.

Physically, this situation models realistic fiber propagation conditions in which the accumulated birefringence-induced phase is not perfectly controlled and varies over time, frequency, or environmental conditions.

The ensemble coherence parameter

\[
\mu
=
\langle e^{i\theta} \rangle
\]

provides a unified physical interpretation of the propagation regime.

Its real part determines the observable transverse correlations and the fidelity with respect to the reference Bell state, while its magnitude quantifies the residual entanglement preserved after averaging.

As the phase fluctuations become broader, the ensemble loses its coherent phase structure. The corresponding density operator progressively approaches a mixed state, leading to a reduction of both purity and concurrence.

This experiment therefore captures the transition from ideal coherent propagation toward realistic noisy optical-fiber evolution, where stochastic birefringence effects induce effective decoherence and entanglement degradation.